

\begin{center}
{\small
\href{mailto:Hkamani@ualberta.ca}{Heena Kamani} \\
\href{mailto:hrho@ualberta.ca}{Hyunsung Rho} \\
\href{mailto:tjongele@ualberta.ca}{Tom Jongeleen}
}
\end{center}

\section{Project Overview}

\subsection{Abstract} 

Real-time crowd simulation for large-scale environments remains a significant challenge in computer graphics due to the high computational cost of collision detection and resolution. While Position-Based Dynamics (PBD) provides a stable and robust framework for multi-agent systems, its performance scales poorly as agent density increases. This project proposes a Computational Level of Detail (C-LOD) framework that dynamically optimizes PBD constraints using Adaptive Spatial Clustering. By intelligently adjusting constraint resolution fidelity based on local agent density and perceptual importance, the proposed system aims to achieve real-time performance for massive crowd simulations without compromising visual quality.

\subsection{1 Introduction} 

The simulation of large crowds has become a cornerstone of modern visual effects and interactive media, while also proving indispensable in robotics, architecture, and industrial safety design. In robotics, crowd modeling informs navigation and multi-agent coordination in dynamic environments. In architectural and industrial contexts, it enables the testing of evacuation routes, pedestrian flow optimization, and the assessment of spatial feasibility in complex infrastructures such as airports and stadiums.  

Regardless of the application, the computational cost of collision detection and resolution represents a significant bottleneck in every crowd simulation. As the number of agents increases, the complexity of interactions grows substantially, leading to exponential increases in computational load. In this project, we propose a novel Computational Level of Detail (C-LOD) framework that dynamically optimizes Position-Based Dynamics (PBD) constraints using Adaptive Spatial Clustering. Our approach aims to reduce the computational cost of collision detection and resolution while maintaining the visual fidelity of the simulation. By clustering agents based on their spatial proximity and dynamically adjusting the level of detail for each cluster, we can significantly improve the performance of crowd simulations without sacrificing visual quality.